\section{An integral interpolation of cyclic characters}
\label{se-section}

The next input prescribes one pushforward of an integral
function on a product of two cyclic groups, while making
several other pushforwards zero.  Working over the integers
is necessary: the function will be used as the exponents
of fractional ideals.  No division by a difference of
slopes is permitted in the group ring.

For a finite abelian group \(A\), write
\(\mathbb Z[A]\) for the free abelian group on symbols
\([a]\), with multiplication \([a][b]=[a+b]\).
Its multiplicative identity is \([0]\), which we also
write as \(1\).
A homomorphism \(f:A\to B\) induces
\(f_*:\mathbb Z[A]\to\mathbb Z[B]\) by
\([a]\mapsto[f(a)]\).
Thus if \(\eta=\sum_a\eta(a)[a]\), the coefficient
of \([b]\) in \(f_*\eta\) is
\(\sum_{f(a)=b}\eta(a)\).

\begin{lemma}\label{se-difference}
Let \(w,M\) be nonnegative integers, put \(\kappa=1+w\),
and write
\[
 P_M(T)=\prod_{j=0}^{M-1}(T-\kappa^j)
                    =\sum_{i=0}^M p_iT^i.
\]
Then in \(\mathbb Z[X]\),
\[
                (X^w-1)^M
                    \ \bigm|\ \sum_{i=0}^M p_iX^{\kappa^i}.
\]
\end{lemma}
\begin{proof}
Let \(U\) be the endomorphism of the abelian group
\(\mathbb Z[X]\) given by \(Uf(X)=f(X^\kappa)\),
and put \(D=X^w-1\).
For every polynomial \(y\), we have
\(Uy\equiv y\pmod D\), since \(X^\kappa\equiv X\pmod D\).
Moreover
\[
 UD=D S,\qquad
 S=\sum_{a=0}^{\kappa-1}X^{wa},\qquad
 S\equiv\kappa\pmod D.
\]
These identities also hold when \(w=0\), with
\(D=0\) and \(U\) the identity.
If \(f=D^jy\), then
\[
 (U-\kappa^j)f
   =D^j\bigl((S^j-\kappa^j)Uy
                      +\kappa^j(Uy-y)\bigr)
\]
is divisible by \(D^{j+1}\).
Starting with \(X\) and applying
\((U-1),(U-\kappa),\ldots,(U-\kappa^{M-1})\)
therefore gives a polynomial divisible by \(D^M\).
The operators commute, and \(U^iX=X^{\kappa^i}\).
Their product applied to \(X\) is consequently the
polynomial in the assertion.
The empty product \(M=0\) gives the same conclusion.
\end{proof}

\begin{lemma}\label{se-divisor}
Let \(\ell\) be prime, \(q=\ell^\nu\), and
\(0\le c\le\nu\).
If \(d\) is a nonzero integer with \(|d|<\ell^c\),
then the image of \(d\) in \(\mathbb Z/q\mathbb Z\)
divides the image of \(\ell^c\).
Consequently, in \(\mathbb Z[\mathbb Z/q\mathbb Z]\),
\[
                  [d]-1\ \bigm|\ [\ell^c]-1.
\]
\end{lemma}
\begin{proof}
Write \(d=\ell^a u\) with \(\ell\nmid u\).
The size hypothesis gives \(a<c\).
The residue class of \(u\) is a unit modulo \(q\),
so a solution to \(dv=\ell^c\) is
\(v=\ell^{c-a}u^{-1}\pmod q\).
Choose a nonnegative integer representative \(v_0\).
Then \([\ell^c]=[d]^{v_0}\), and the usual geometric
sum proves divisibility by \([d]-1\).
If \(v_0=0\), the numerator is zero and divisibility
still holds.
\end{proof}

\begin{theorem}\label{se-interpolation}
Let \(\ell\) be prime and let \(c,\nu,M\) be
nonnegative integers satisfying
\[
                 c\le\nu,\qquad M<\ell^c.
\]
Put \(q=\ell^\nu\), \(A=\mathbb Z/q\mathbb Z\),
and \(\kappa=1+\ell^c\).
For \(0\le t\le M\), define
\[
             \chi_t:A^2\longrightarrow A,
             \qquad \chi_t(x,y)=x+ty.
\]
Fix \(t_0\in\{0,\ldots,M\}\), and write
\[
              P_M(T)=\prod_{j=0}^{M-1}(T-\kappa^j)
                       =\sum_{i=0}^Mp_iT^i.
\]
There exists \(\eta\in\mathbb Z[A^2]\) such that
\[
  (\chi_t)_*\eta=0\quad(t\ne t_0),
  \qquad
  (\chi_{t_0})_*\eta
                       =\sum_{i=0}^M p_i[\kappa^i].
\]
\end{theorem}
\begin{proof}
All residue classes in this proof lie in \(A\).
Let
\[
          \Delta=\sum_{i=0}^Mp_i[\kappa^i],
          \qquad D=[\ell^c]-1.
\]
Apply the ring homomorphism
\(\mathbb Z[X]\to\mathbb Z[A]\), \(X\mapsto[1]\),
to Lemma~\ref{se-difference}.
It gives \(D^M\mid\Delta\).
For each \(t\ne t_0\) in the indicated range,
\(0<|t_0-t|\le M<\ell^c\).
Lemma~\ref{se-divisor} therefore gives
\([t_0-t]-1\mid D\).
There are exactly \(M\) such slopes, so
\[
          B:=\prod_{t\ne t_0}([t_0-t]-1)
                      \ \bigm|\ D^M\ \bigm|\ \Delta.
\]
Choose \(Q\in\mathbb Z[A]\) with \(BQ=\Delta\).
This argument multiplies divisibility identities;
it never cancels a zero divisor in the group ring.

Let \(\iota:\mathbb Z[A]\to\mathbb Z[A^2]\)
be induced by \(a\mapsto(a,0)\), and put
\[
 \eta=
 \prod_{t\ne t_0}\bigl([(0,1)]-[(t,0)]\bigr)
 \ \iota\!\left(\left[-\sum_{t\ne t_0}t\right]Q\right).
\]
If \(s\ne t_0\), the factor indexed by \(t=s\)
has zero image under \((\chi_s)_*\), so
\((\chi_s)_*\eta=0\).
For the distinguished slope,
\[
 (\chi_{t_0})_*\bigl([(0,1)]-[(t,0)]\bigr)
       =[t_0]-[t]=[t]([t_0-t]-1).
\]
Also \((\chi_{t_0})_*\iota\) is the identity.
The product of the factors \([t]\) cancels the
explicit unit \([-\sum t]\), leaving \(BQ=\Delta\).
This proves both required pushforward identities.
\end{proof}

In the arithmetic cycle construction,
\(M=2E+1\), and the hypothesis is
\(2E+1<\ell^c\).
The proof retains the exact integer coefficients, so
they can be used as positive or negative powers of
fractional ideals.
It also includes \(M=0\), when the product is empty,
and \(c=\nu\), when \([\ell^c]-1=0\).
In the latter case with \(M>0\), the target
\(\Delta\) is zero by the proved divisibility;
there is no illicit division by zero.


